\input{../tex-inputs/latex-leseansicht-vorspann}

\section[Arthur und Olga Schnitzler an Gustav Schwarzkopf, 5. 5. 1914]{L04492 Arthur und Olga Schnitzler an Gustav Schwarzkopf, 5. 5. 1914}
\nopagebreak\mylabel{L04492v}
\rehead{ }\normalsize\beginnumbering\briefempfaengerindex{Schwarzkopf, Gustav@\textsc{Schwarzkopf, Gustav}!zzzSchnitzler, Olga@\emph{von Olga Schnitzler}!1914-05-052@{5. 5. 1914}|(be}\briefempfaengerindex{Schwarzkopf, Gustav@\textsc{Schwarzkopf, Gustav}!zzzSchnitzler, Arthur@\emph{von Arthur Schnitzler}!1914-05-052@{5. 5. 1914}|(be}
\toendnotes[C]{\smallbreak\pagebreak[2]}
\correspDesc{Versand  durch Arthur Schnitzler, Olga Schnitzler am 5. 5. 1914 in Florenz
\newline{}Übermittlung  am 5. 5. 1914 in Florenz
\newline{}Erhalt  durch Gustav Schwarzkopf im Zeitraum [6. 5. 1914 – 10. 5. 1914?] in Wien}\toendnotes[C]{\smallbreak}
\Standort{DLA, A:Schnitzler, HS.1985.1.1897.}
\physDesc{Bildpostkarte, 211 Zeichen
\newline{}Handschrift Arthur Schnitzler: Bleistift, deutsche Kurrent
\newline{}Handschrift Olga Schnitzler: Bleistift, lateinische Kurrent
\newline{}Versand: Stempel: »\nobreak{}\oindex{Florenz@\textbf{Florenz}|pwk}Firenze, 5. V. 1914, 20–21, ferrovia\nobreak{}«.  }\toendnotes[C]{\smallbreak}\pstart{}{\pb}Hrn \textsc{Gustav Schwarzkopf}\pend{}\pstart{}\textsc{Wien I}\oindex{I., Innere Stadt@\textbf{I., Innere Stadt}|pw}\pend{}\pstart{}\textsc{Tiefer Graben 17}\oindex{Wien@\textbf{Wien}!I., Innere Stadt@\textbf{I., Innere Stadt}!Tiefer Graben 17@\textbf{Tiefer Graben 17}|pw}\pend{}{\bigskip}
\pstart
           \noindent{}{\pb}\textcolor{gray}{\textbf{Firenze – Piazza della Signoria\oindex{piazza della Signoria@\textbf{piazza della Signoria}|pw} – Fontana del Nettuno\pwindex{Neptunbrunnen (Florenz)@\emph{Neptunbrunnen (Florenz)}|pw}}}{ }\textcolor{gray}{\textbf{(B. Ammannati\pwindex{Ammannati, Bartolomeo 18.\,6.\,1511 Settignano – 13.\,4.\,1592 Florenz@\textsc{Ammannati, Bartolomeo} (18.\,6.\,1511 Settignano – 13.\,4.\,1592 Florenz)|pw})}}\pend
           \vspace{1em}
\pstart
           \raggedleft{}{\pb}Florenz\oindex{Florenz@\textbf{Florenz}|pw}{ }5. 5. 914\pend
           \vspace{0.5em}
\pstart
           Herzliche Grüße! Bisher
      alles ſo ſchön und angenehm als möglich. Möge
      es ſo bleiben. Warum{ }ſoll man nicht{ }ſich{ }ſelber
      was wünſchen?\pend
           \pstart Ihr \spacefill\mbox{A.}\pend{}\selectlanguage{ngerman}\vspace{1em}
\pstart
           \noindent{}{[}hs. O. Schnitzler:{]} \label{T_L04492-1v}\edtext{Herzlichst \spacefill\mbox{Olga}}{\lemma{\textnormal{\emph{Herzlichst Olga}}}\Cendnote{\textnormal{Der Gruß befindet sich entlang des linken Randes der Karte.}}}\label{T_L04492-1}\pend
           \selectlanguage{ngerman}\endnumbering\briefempfaengerindex{Schwarzkopf, Gustav@\textsc{Schwarzkopf, Gustav}!zzzSchnitzler, Olga@\emph{von Olga Schnitzler}!1914-05-052@{5. 5. 1914}|)be}\briefempfaengerindex{Schwarzkopf, Gustav@\textsc{Schwarzkopf, Gustav}!zzzSchnitzler, Arthur@\emph{von Arthur Schnitzler}!1914-05-052@{5. 5. 1914}|)be}\mylabel{L04492h}
\begin{anhang}
\end{anhang}\newcommand{\dateiname}{L04492}\newcommand{\titel}{Arthur und Olga Schnitzler an Gustav Schwarzkopf, 5. 5. 1914}\newcommand{\editorInnen}{Herausgegeben von Jahnke, SelmaMüller, Martin Anton}\input{../tex-inputs/latex-leseansicht-abspann}
